\subsection{Ведущий и ведомый}
\label{sec:swarm-follow}

Цель раздела: использовать телеметрию соседа в собственном полёте --- лететь за ведущим на заданном расстоянии.

Идея: ведущий просто выполняет свой маршрут и ни о чём не знает; ведомый раз в полсекунды читает позицию ведущего из его телеметрии и летит в точку со смещением. Ведущему \emph{не нужно ничего отправлять}: данные берутся из публичной телеметрии. Два файла: ведомый запускается первым, ведущий --- вторым.

\begin{codefile}[python]{python-onboard/swarm/swarm_05_follow_follower.py}
\end{codefile}

\paragraph{Что происходит (ведомый).} Идентификатор ведущего \verb|LEADER = 'pioneer1'| задан константой; смещение \verb|OFFSET_X = -1.0| означает «на метр позади по оси $x$». Шестьдесят раз с шагом 0.5~с ведомый читает позицию ведущего (\verb|get_telemetry_from_unit| с тайм-аутом 2~с) и отправляет \verb|go_to_local_point| в точку «позиция ведущего плюс смещение» на высоте 1~м. Если ведущий пропал из сети (\verb|None|), цикл прекращается; после цикла дрон садится. Обратите внимание: ведомый вызывает команду точки, не дожидаясь прибытия, --- следующая точка просто заменяет предыдущую.

\begin{codefile}[python]{python-onboard/swarm/swarm_05_follow_leader.py}
\end{codefile}

\paragraph{Что происходит (ведущий).} Ведущий взлетает и облетает квадрат $2\times 2$~м: четыре точки с ожиданием \verb|point_reached()| в каждой и секундной паузой. Затем садится. Работа с \verb|swarm| ограничена созданием и закрытием: ведущий присутствует в сети уже потому, что у его борта работает \texttt{swarm-server}.

\begin{warnbox}[Что нужно, чтобы следование работало на настоящих дронах]
Ведомый подставляет координаты ведущего в \emph{свою} систему координат. Это верно, только если у обоих \emph{общая} система --- например, оба летят по одной карте локальной системы навигации или стартуют из одной точки (так в симуляторе после сброса). Если у каждого дрона своё начало координат, следование приведёт ведомого не туда. Кроме того, в примере нет защиты от столкновений, нет контроля курса и запаса по батарее. Проверьте общие координаты на земле: поставьте оба дрона в разные точки и сравните их локальные позиции.
\end{warnbox}

\begin{taskbox}[Ведомый и ведущий]
Запустите ведомого на втором дроне, затем ведущего на первом. Измените смещение ведомого на $1$~м слева и $1$~м позади и убедитесь по картинке 3D, что он держится ближе к ведущему слева. Добавьте в ведомый печать расстояния до ведущего.
\expected{ведущий печатает четыре строки «Ведущий в точке \ldots»; ведомый летит следом с отставанием в полсекунды; оба садятся; расстояние до ведущего колеблется около выбранного смещения (порядка 1.4~м при смещении по обеим осям).}
\mission{Мини~2 SDK2}{62}
\end{taskbox}

\subsection{Строй: у каждого своё место}
\label{sec:swarm-formation}

Один ведущий и несколько ведомых образуют строй: каждому ведомому нужно \emph{своё} место --- смещение от ведущего. Одну и ту же программу запускают на нескольких дронах, а место выбирают по \texttt{machine\_id}.

\begin{codefile}[python]{python-onboard/swarm/swarm_08_formation_follower.py}
\end{codefile}

\paragraph{Что происходит.} Список \verb|SLOTS| --- места: $(-1,\,1)$, $(-1,\,-1)$, $(-2,\,0)$ (метры в осях сцены). Номер дрона берётся из идентификатора (\verb|int(swarm.id.replace('pioneer', ''))|), и место --- \verb|SLOTS[(number - 2) % len(SLOTS)]|: \texttt{pioneer2} занимает первое, \texttt{pioneer3} --- второе. Остальное --- как в предыдущем примере. Остаток от деления гарантирует, что при большем числе дронов места начнут повторяться, а не вызовут ошибку; в реальной программе это нужно обрабатывать отдельно: два дрона на одном месте столкнутся.

Ведущего для строя запускают так же, как в прошлом разделе (\texttt{swarm\_05\_follow\_leader.py}). Раздача мест по номеру в имени проста, но хрупка: она рассчитана на \texttt{pioneer2} и далее. Надёжнее договориться сообщениями: ведущий рассылает каждому его место.

\begin{taskbox}[Треугольник]
Запустите ведомых на втором и третьем дронах, ведущего --- на первом. Затем замените раздачу мест по номеру на раздачу сообщением: ведущий перед взлётом отправляет каждому ведомому строку \verb|'место dx dy'|, а ведомый разбирает её и только потом взлетает.
\expected{в 3D два ведомых летят справа и слева позади ведущего; при отсутствии сообщения ведомый не взлетает и печатает причину.}
\end{taskbox}

\subsection{Синхронный взлёт по готовности}

Группа не должна взлетать, пока кто-то не готов. Решение --- протокол из трёх сообщений: командир посылает «готовность», каждый ведомый запускает моторы и отвечает «готов», командир ждёт всех и посылает «взлёт».

\begin{codefile}[python]{python-onboard/swarm/swarm_06_sync_takeoff_wingman.py}
\end{codefile}

\begin{codefile}[python]{python-onboard/swarm/swarm_06_sync_takeoff_commander.py}
\end{codefile}

\paragraph{Что происходит.} Ведомый (запускается раньше) ждёт «готовность» функцией \verb|wait_for(command, limit)|: опрос ящика с паузой и пределом, который возвращает отправителя, или \verb|None|, если пришло «отбой» или время вышло. Получив «готовность», ведомый запускает моторы (\verb|arm()|) и отвечает «готов». Затем ждёт «взлёт» и взлетает, а по «посадка» садится; без «взлёта» выключает моторы. Командир собирает ответы в множество \verb|ready| (не дольше 20~секунд), и только когда готовы \emph{все}, делает обратный отсчёт, рассылает «взлёт» и взлетает сам; иначе рассылает «отбой» и ничего не запускает. Моторы запущены до отсчёта, поэтому после «взлёта» дроны поднимаются почти одновременно.

\begin{rulebox}[Не взлетай, пока все не ответили]
Протокол группы строят вокруг подтверждений: «готовность» $\to$ «готов» $\to$ «взлёт». Общее действие начинается только после того, как подтвердили все, а отказ любого участника (тишина, «отбой») отменяет всё. Предел ожидания нужен всегда: иначе один неотвечающий дрон держит группу навсегда.
\end{rulebox}

\begin{taskbox}[Взлёт по готовности]
Запустите ведомых на втором и третьем дронах, затем командира. Посмотрите, что произойдёт, если не запускать одного из ведомых: сколько ждёт командир и что напечатает.
\expected{при двух ведомых командир печатает «pioneer2 готов», «pioneer3 готов», отсчёт 3, 2, 1 и «Все в воздухе»; все три дрона взлетают и садятся; если ведомый не запущен, через 20~с командир печатает «Не все готовы, взлёт отменён».}
\mission{Мини~2 SDK2}{63}
\end{taskbox}

\subsection{Световое шоу}

Тот же приём --- команды-слова --- подходит для синхронной подсветки: дирижёр рассылает название цвета, а каждый игрок включает его.

\begin{codefile}[python]{python-onboard/swarm/swarm_07_led_show_player.py}
\end{codefile}

\begin{codefile}[python]{python-onboard/swarm/swarm_07_led_show_conductor.py}
\end{codefile}

\paragraph{Что происходит.} Игрок хранит словарь \verb|COLOURS| (название $\to$ тройка 0..1) и опрашивает ящик с паузой 0.05~с; незнакомое слово игнорируется (мусор не роняет программу), \verb|'off'| завершает шоу. Дирижёр дважды проходит цвета красный, зелёный, синий, белый: рассылает название каждому соседу, включает тот же цвет себе (\verb|led_control(255, r, g, b)| --- все светодиоды) и ждёт секунду. В конце рассылает \verb|'off'| и гасит свои светодиоды. Рассылку по списку соседей при каждом цвете делают заново: за время шоу состав мог измениться.

\begin{taskbox}[Свой номер шоу]
Придумайте программу шоу: бегущий цвет по дронам (в каждый момент светится один), затем все вместе, затем гаснут по очереди. Для этого дирижёр рассылает не одно слово, а пару «номер и цвет».
\expected{в журнале игроков команды появляются в порядке сценария; финальное сообщение гасит светодиоды у всех участников, включая дирижёра.}
\end{taskbox}
